\section{Introduction}
\label{sec:intro}

Matrix multiplication is a core computational primitive underlying many large-scale computations, including numerical linear algebra, signal processing, computer graphics, and artificial intelligence.  Consequently, many applications require verifiability for matrix products while keeping the underlying matrices private.  Zero-knowledge proofs provide a cryptographic mechanism for this goal:
a prover can convince a verifier that a statement is correct without revealing any secret information that witnesses the statement.

Applying general-purpose ZKP systems directly to large matrix multiplication, however, remains expensive.  Succinct proof systems such as ~\cite{groth2016size, plonk, maller2019sonic, chiesa2020marlin} prove relations after the computation has been represented as an arithmetic circuit or a similar algebraic constraint system.  
A standard arithmetization of a \(k\times k\) matrix product \(\mat{A}\mat{B}=\mat{C}\) computes every output
entry and therefore uses \(\mathcal{O}(k^3)\) arithmetic constraints.  This is
already impractical for moderate \(k\), and it is especially problematic for
Transformer-based neural networks~\cite{transformer, bert, GPT3}, whose projection and attention layers often
contain matrix multiplications with dimension \(k\ge 1{,}000\). \\

\noindent\textbf{The \(\mathcal{O}(k^2)\) barrier.}
Freivalds' classical randomized check~\cite{freivalds1977probabilistic} reduces
matrix-product verification from a cubic computation to a quadratic one.  Given
matrices \(\mat{A},\mat{B},\mat{C}\) and a random vector \(\vect{r}\), the
verifier checks
\[
  \vect{r}\mat{A}\mat{B}=\vect{r}\mat{C}
\]
instead of recomputing the full product.  This reduces the verification cost
from \(\mathcal{O}(k^3)\) to \(\mathcal{O}(k^2)\).  The resulting
\(\mathcal{O}(k^2)\)-size circuit is an improvement, but it is still expensive:
for example, at \(k=2^{12}\), it already requires \(50.5\) million R1CS
constraints and \revmath{411.28\,\text{s}} of proving time with Groth16.  Thus, a large gap remains
between the \(\mathcal{O}(k^2)\) SNARK cost and practical large-scale
matrix-multiplication verification. \\

\noindent\textbf{Our insight: from computation to checking.}
We observe that the \(\mathcal{O}(k^2)\) barrier is not inherent to the
verification problem itself.  It arises because existing SNARK-based approaches
still \emph{compute} the vector--matrix products inside the circuit.  Our key
idea is to move from ``computing inside the circuit'' to ``\emph{checking} inside the
circuit.''  The prover performs the expensive matrix arithmetic outside the
SNARK and commits to the relevant encoded data and intermediate values.  The
verifier then checks only a small number of lightweight algebraic relations
that are sufficient to detect inconsistent claims.
The technical core of this approach is proximity testing over linear
error-correcting codes.  Let \(\mathcal{C}\) be a linear code with relative
distance \(\delta\).  We encode each matrix row-wise and consider the Freivalds
intermediate vectors
\[
  \vect{x}=\vect{r}\mat{A},\qquad
  \vect{y}=\vect{x}\mat{B},\qquad
  \vect{z}=\vect{r}\mat{C}.
\]
The linearity of \(\mathcal{C}\) lets us check these vector--matrix products in
the encoded domain, column by column.  For one encoded column, a single inner
product, or \emph{fold}, checks one coordinate of the encoded product.  If the
prover cheats in one of the corresponding relations, the two relevant
codewords disagree on at least a \(\delta\)-fraction of positions.  A random
query therefore detects the inconsistency with probability at least \(\delta\).

The key point is that the SNARK circuit never takes the full matrices as witnesses.  Instead, it checks fold equations on \(t\) sampled encoded columns, where \(t\) is independent of the matrix dimension \(k\).  It also checks that the committed encoded views of the intermediate vectors \(\vect{x},\vect{y},\vect{z}\) are valid codewords for those vectors.  
To bind the circuit witnesses to the data fixed before the query positions are known,
we use a commit-and-prove SNARK for the sampled block witnesses.  The prover
also supplies Merkle membership proofs for the queried leaves under
\(\rt_{ABC}\) and \(\rt_{XYZ}\), together with CP-Link proofs showing that the
SNARK-side sampled blocks and the opened Pedersen leaves contain the same data.
These checks ensure, under the backend assumptions, that the values checked
inside the circuit are consistent with commitments fixed before the verifier's
sampling challenges.

% As a result, the main SNARK relation has \(\mathcal{O}(tk)\) constraints.
% For fixed \(t\), this is linear in the matrix dimension.  The sampled opening
% and linking layer adds \(\mathcal{O}(t\log n)+E_{\mathsf{link}}(k,t)\) backend
% work.  In our implementation, verification consists of \(\mathcal{O}(1)\) pairings for SNARK and CP-Link verification and \(\mathcal{O}(t\log n)\) hash operations for Merkle openings.\\

As a result, the main SNARK relation has
\(\mathcal{O}(tk+E_{\mathsf{code}})\) constraints, where
\(E_{\mathsf{code}}\) is the cost of checking that the
intermediate vectors are valid encodings under the error-correcting code. For the fixed-rate Reed--Solomon setting used in our implementation,
this remains linear in the matrix dimension for fixed \(t\).  The sampled
opening and linking layer adds
\(\mathcal{O}(t\log n)+E_{\mathsf{link}}(k,t)\) backend work.  In our
implementation, verification consists of SNARK verification, CP-Link
verification, and \(\mathcal{O}(t\log n)\) hash operations for Merkle openings;
for the QALink instantiation, CP-Link verification includes
\(\mathcal{O}(t)\) pairing terms batched into a multi-pairing check.\\


\noindent\textbf{Contributions.}
We present \(\protocol\), a protocol for efficiently proving
matrix-multiplication statements.  Our contributions are as follows.
\begin{itemize}
  \item \textbf{A code-based matrix-product checking protocol.}
  \(\protocol\) combines Freivalds' randomized reduction with proximity testing
  over linear error-correcting codes.  Instead of computing vector--matrix
  products inside the SNARK circuit, the protocol checks sampled fold equations
  over encoded commitments and verifies code consistency for the intermediate
  vectors. 

  For fixed \(t\), the main circuit cost is
  \(\mathcal{O}(tk+E_{\mathsf{code}})\), which is linear in \(k\) for the
  fixed-rate Reed--Solomon setting used in our implementation.

  \item \textbf{Soundness analysis.}
  \begin{revblock}
  We prove soundness in the committed-input setting by leveraging the relative
  distance and unique-decoding radius of the Reed--Solomon code.  The analysis
  also accounts for Merkle openings and CP-Link proofs, which bind the values
  checked inside the SNARK circuit to the data committed before the sampling
  challenges are derived.
  \end{revblock}

  \item \textbf{Implementation and evaluation.}
  \begin{revblock}
  We implement \(\protocol\) in Go and evaluate it against a Freivalds-based
  SNARK baseline and DualMatrix.  In the square-matrix experiments with
  \(\rho=1/2\) and \(t=309\), \(\protocol\) reduces the constraint count by
  \(9.68\times\) and proving time by \revmath{5.81\times} relative to the Freivalds
  baseline at \(k=2^{12}\).  It also proves faster than DualMatrix at every
  matrix dimension measured for both systems, achieving a \revmath{3.62\times}
  speedup at \(k=2^{13}\).

  We further evaluate batching and a GPT-2 workload.  Batching matrix products
  with a shared leaf layout keeps the Merkle proof size and verification time
  nearly constant as the number of claims increases.  On GPT-2, \(\protocol\)
  proves faster than both the Freivalds baseline and DualMatrix, although
  heterogeneous matrix dimensions require multiple commitment layouts and
  increase proof size and verification cost.
  \end{revblock}
\end{itemize}

% \noindent\textbf{Paper organization.}
% Section~\ref{sec:related} discusses related work.
% Section~\ref{sec:prelim} introduces the necessary background.
% Section~\ref{sec:tech-overview} gives an overview of \(\protocol\), and
% Section~\ref{sec:construction} presents the full construction.
% Section~\ref{sec:security} proves security, and
% Section~\ref{sec:implementation} reports the implementation and evaluation.
% Section~\ref{sec:conclusion} concludes.
% Additional proofs and backend details appear in the appendices.
